\section{Introduction}
A service provider can improve its operating policy and still be unable to offer it. Customers with renewal or cancellation rights evaluate what the agreement promises from the current review onward. The provider must therefore respect participation at each continuation, not only the surplus forecast when the agreement was signed. A second constraint is organizational: a service rule may use a preannounced schedule, the current operating regime, a small memory record, or the entire observed history. The value of adaptation must compare the best accepted contracts under the competing information architectures.

This paper identifies a class in which the information needed for an optimal accepted contract can be constructed exactly without enumerating histories. The class has public exogenous Markov uncertainty, continuous service tiers, quadratic operating rewards, positive additive payments, and continuation payment caps. For the principal computational theorem, changing a tier carries no intertemporal switching cost. Participation constraints may bind at any review; their cumulative effect is the central feature, not a small perturbation.

The key object is the payment response to a continuation price. A local service decision has two tier thresholds. At each public state, imposing its continuation cap adds at most one further threshold. Backward propagation therefore produces a finite global set of price breakpoints whose size grows linearly with the public graph, rather than with its exponentially larger history tree. The provider executes the optimizer by retaining the largest participation-price threshold encountered so far. This gives an explicit exact algorithm, a bounded state representation, and independently auditable primal--dual equality. No full-tree optimum or full-tree dual solution is assumed available to generate the representation.

The price record has operational content. Two customers reaching the same public regime may require different service tiers because earlier renewal constraints imposed different effective prices. We prove that the worst-case number of distinguishable price labels grows linearly with the public graph. For a family of renewal agreements, we derive the exact value attainable with a limited number of memory symbols. The optimal encoding groups ordered continuation caps into contiguous cells, and a finite dynamic program computes the resulting value--memory frontier. This specifies what a memory device produces before any installation cost is assigned to it.

A shared policy table remains a global commitment. We derive a graph-sized quadratic program for its optimized zero-release value, including an extension with switching costs. At positive release, the finite-price algorithm supplies exact fixed-table values and supporting prices for a Benders-style shared-design master. The outer master is standard, and no universal finite exact cut count is inferred. The constructive finite-price result is distinguished throughout from the more general extensive-form-dual representation retained for switching and promise-state models.

The paper keeps the broader mathematical framework rather than replacing it with the tractable subclass. Print appendices retain the release frontier, comparator-adjusted continuation rents, friction paths, general promise recursion, shared-table representation, and deployment certificates. Their complete supporting proofs remain in the electronic companion. These results describe accepted adaptation beyond the principal finite-price theorem; they are not additional claims of new decomposition principles.

Exact experiments stress horizon length, public-state recombination, binding continuations, fixed shared corridors, and optimized tables. They include independent unfolded-tree matrix checks and a separate numerical optimizer on small instances. No field data or calibration are claimed. The earlier adverse computational findings remain visible: in the inherited strict-tolerance cache experiment, 83 of 84 queries refreshed and the cached pipeline was not faster. That experiment is not repurposed as evidence for the new algorithm. The new contribution is the constructive finite-price theorem and its memory consequences, supported by exact coupled-graph computations.
